\documentclass[USenglish]{article}	
% for 2-column layout use \documentclass[USenglish,twocolumn]{article}

\usepackage[utf8]{inputenc}				%(only for the pdftex engine)
%\RequirePackage[no-math]{fontspec}[2017/03/31]%(only for the luatex or the xetex engine)
\usepackage[big,online]{dgruyter}	%values: small,big | online,print,work
\usepackage{lmodern} 
\usepackage{microtype}
\usepackage[numbers,square,sort&compress]{natbib}
\usepackage{algorithm}
\usepackage[utf8]{inputenc} % allow utf-8 input
% babel is loaded by dgruyter.sty; keep English only to avoid option clash
% \usepackage[english]{babel}
%\usepackage[T1]{fontenc}    % use 8-bit T1 fonts
\usepackage{booktabs}       % professional-quality tables
\usepackage{amsfonts}       % blackboard math symbols
\usepackage{nicefrac}       % compact symbols for 1/2, etc.
\usepackage{microtype}      % microtypography
\usepackage{lipsum}		% Can be removed after putting your text content
\usepackage{graphicx}
\usepackage{natbib}
\usepackage{doi}
\usepackage{amsmath}
\usepackage{amsthm}
\usepackage{amssymb}
\usepackage{mathtools}
\usepackage{algorithm}
\usepackage{algpseudocode}
% dgruyter.sty hooks caption and expects \plist@<float-type> for every
% float family that writes into the corresponding *.lo? file. The
% algorithm package writes into .loa, so we provide the missing macro.
\makeatletter
\providecommand\plist@algorithm{\algorithmcfname\space}
\providecommand\algorithmcfname{Algorithm}
\makeatother
\usepackage{appendix}

\usepackage{caption}
\usepackage{subcaption}
\usepackage{wrapfig}


\graphicspath{{logos/}, {images/}, {}}
\DeclareGraphicsExtensions{.pdf,.png,.jpg}

% Цвета 
\usepackage[dvipsnames]{xcolor}
\usepackage{color}   


% New theorem-like environments will be introduced by using the commands \theoremstyle and \newtheorem.
% Please note that the environments proof and definition are already defined within dgryuter.sty.
\theoremstyle{dgthm}
\newtheorem{theorem}{Theorem}
\newtheorem{assumption}{Assumption}
\newtheorem{corollary}{Corollary}
\newtheorem{proposition}{Proposition}
\newtheorem{lemma}{Lemma}
\newtheorem{assertion}{Assertion}
\newtheorem{result}{Result}
\newtheorem{conclusion}{Conclusion}

\newtheorem{definition}{Definition}
\newtheorem{example}{Example}
\newtheorem{remark}{Remark}



\begin{document}

	
%%%--------------------------------------------%%%
	\articletype{Research Article}
	\received{Month	DD, YYYY}
	\revised{Month	DD, YYYY}
  \accepted{Month	DD, YYYY}
  \journalname{De~Gruyter~Journal}
  \journalyear{YYYY}
  \journalvolume{XX}
  \journalissue{X}
  \startpage{1}
  \aop
  \DOI{10.1515/sample-YYYY-XXXX}
%%%--------------------------------------------%%%

\title{Structure Aware Neural Architecture Search for Mixture of Experts}
\runningtitle{NAS for MoE}
%\subtitle{Insert subtitle if needed}

\author[1]{Petr Babkin}
\author[1]{Oleg Bakhteev}

\affil[1]{\protect\raggedright 
Moscow Institute of Physics and Technology, Intelligent Systems, Moscow, Russia, e-mail: babkin.pk@phystech.edu}

%\communicated{...}
%\dedication{...}
	
\abstract{The Mixture-of-Experts (MoE) layer, a sparsely activated neural architecture controlled by a routing mechanism, has recently achieved remarkable success across large-scale deep learning tasks. In parallel, Neural Architecture Search (NAS) has emerged as a powerful methodology for automatically discovering high-performing neural network. However, the application of NAS methods to MoE architectures remains an underexplored research area. In this work, we propose an architecture search framework for MoE models, which explicitly leverages the underlying cluster structure of the data. We evaluate the proposed approach on computer vision benchmarks and demonstrate that it outperforms baseline MoE architectures trained on the same datasets in terms of accuracy and computational efficiency.}

\keywords{Mixture of Experts, Neural Architecture Search}

\maketitle

\section{Introduction}
\label{sec:introduction}


The Mixture-of-Experts (MoE) architecture, in which input is partitioned through a gating function and subsequently routed as weighted input to specialized experts \cite{jacobs1991adaptive, jordan1994hierarchical, jacobs1997bayesian}, has been well-established in the literature for decades \cite{mu2025comprehensive}.
Scenarios in which the experts are classical statistical models are well characterized in the literature \cite{piwko2025divide}, particularly when the experts are modeled as logistic regressors \citep{nguyen2023general} or as Gaussian models \citep{nguyen2023demystifying, rasmussen2001infinite}.
In recent years, concurrent with the growing interest in neural networks, researchers began exploring the application of MoE structures to deep learning architectures \cite{eigen2013learning}.
This efficiency was achieved through sparse gating mechanisms, where the input is routed only to a fixed number of experts, which is significantly smaller than the complete set \cite{shazeer2017outrageously, riquelme2021scaling, fedus2022switch}.
MoE architectures found widespread adoption in the transformer domain, where model size plays a crucial role in achieving superior performance \cite{jiang2024mixtral, qiu2025demons, dai2024deepseekmoe}.

Simultaneously, the field of Neural Architecture Search (NAS) has developed into a well-studied area \cite{ren2021comprehensive, nayman2019xnas, zoph2018learning}.
NAS aims to select optimal architectures from a wide range of candidates, addressing the challenge that architecture selection typically requires substantial domain expertise and manual tuning \cite{liu2018darts, ren2021comprehensive}.
The field has diverged into two main paradigms.
The first is computationally intensive many-shot NAS, which involves multiple training runs followed by candidate selection using methods such as reinforcement learning \cite{zoph2016neural, baker2016designing} and evolutionary algorithms \cite{real2019regularized, yang2020cars}.
The second is one-shot NAS, where a single supernet is trained, and the optimal subnet is subsequently selected \cite{liu2018darts, chen2020stabilizing, you2020greedynas}.
This approach effectively transforms the search problem into finding optimal compression strategies for the supernet \cite{liu2018darts}.

Despite the individual successes of both MoE and NAS methodologies, the application of NAS techniques specifically to MoE architectures remains a relatively underexplored research area, presenting significant opportunities for advancing both efficiency and performance in large-scale neural architectures.
In contrast to the classical approach where all experts share the same architecture \citep{fedus2022switch}, our study allows architectural heterogeneity between experts.
We claim that this flexibility enhances the overall efficiency of the resulting architecture by enabling more specialized and adaptable expert networks.

It is important to note that several studies in the literature have addressed the problem of discovering or designing architectures for mixture-of-experts models. In the article \cite{jawahar2022automoe}, the authors employed a comprehensive combination of approaches to discover a fast-operating Mixture-of-Experts (MoE) architecture. The expert weights are derived from a supernet, and the search process is evolutionary in nature. There is a practical work \cite{han2024cmn} in which the authors develop a framework for efficient MoE inference tailored to specific hardware. The work most closely related to ours is \cite{mecharbat2025moenas}, where the authors conduct a search by varying the number of experts in the layers. In contrast to this study, we propose to allow greater flexibility by optimizing the model architectures themselves rather than only the number of experts.

The MoE architecture is inherently well-suited for data with strong clustering properties. There exist theoretical studies, which demonstrate on relatively simple examples that this claim holds \cite{chen2022towards}.
Additionally, practical researches show the effectiveness of MoE on multimodal datasets \cite{chai2022mixture, guo2018multi, hu2025multi}.
However, despite these promising results, the underlying mechanisms that enable experts to align with distinct modalities remain insufficiently understood, motivating further investigation into the principles governing modality-specific specialization within MoE frameworks.

Given the foregoing, we formulate the following research question:
\textit{how can we design conditions under which expert architectures adapt to specific data clusters?}
In this study, we aim to address this question by examining how models evolve during adaptation to specific data modalities.









\vspace{2mm}

\textbf{Contributions}

\vspace{2mm}

\begin{itemize}
    \item \textbf{Problem formulation.} We propose a new formulation of architecture search for Mixture-of-Experts models as a cluster-aware likelihood-maximisation problem, and give a theoretical grounding for it by showing its equivalence to the incomplete-data maximum likelihood of a latent-variable mixture.
    \item \textbf{Algorithm.} We propose an algorithm that solves the stated problem, based on the Expectation--Maximisation (EM) algorithm with an adaptively refined surrogate model.
    \item \textbf{Convergence guarantee.} We prove a convergence theorem for the proposed algorithm.
    \item \textbf{Benchmarks and experiments.} We construct datasets on which mixture-of-experts approaches have a clear advantage, and conduct comparative experiments against several baselines on them.
\end{itemize}

\section{Related Work}
\label{sec:related_work}

\noindent

\vspace{2mm}

\textbf{Mixture-of-Experts.}
The Mixture-of-Experts (MoE) framework was originally introduced
in~\cite{jacobs1991adaptive,jordan1994hierarchical,jacobs1997bayesian}
as a way to combine specialised statistical models through a learnable
gating mechanism. Subsequent theoretical analyses characterised
identifiability and convergence in classical MoE
formulations~\cite{nguyen2023general,nguyen2023demystifying,piwko2025divide,rasmussen2001infinite}.
With the rise of deep learning, MoE has been revived in the form of
sparsely activated neural networks, where the gate routes each input
to only a small subset of
experts~\cite{eigen2013learning,shazeer2017outrageously,riquelme2021scaling,fedus2022switch}.
This sparse formulation enables training of models with hundreds of
billions of parameters while keeping per-token compute roughly
constant, and is now a standard building block of modern large-scale
transformers~\cite{jiang2024mixtral,dai2024deepseekmoe,qiu2025demons}.

\vspace{2mm}

\textbf{Neural Architecture Search.}
Neural Architecture Search (NAS) automates the design of neural
networks and has matured into a wide research
field~\cite{ren2021comprehensive}. Many-shot approaches based on
reinforcement learning~\cite{zoph2016neural,baker2016designing,zoph2018learning}
or evolutionary search~\cite{real2019regularized,yang2020cars}
explicitly train candidate architectures and select the best
performer. One-shot methods reduce the search cost by training a
single supernet that contains all candidate architectures as
sub-networks~\cite{liu2018darts,chen2020stabilizing,you2020greedynas,nayman2019xnas}.
DARTS~\cite{liu2018darts} is the most prominent representative of this
family and forms one of the baselines we compare against in our
experimental study.

\vspace{2mm}

\textbf{Neural Architecture Search for Mixture-of-Experts.}
The intersection of NAS and MoE is a comparatively recent and still
sparsely populated research area. AutoMoE~\cite{jawahar2022automoe}
applies an evolutionary search on top of a shared supernet to obtain
heterogeneous MoE architectures for neural machine translation.
CMN~\cite{han2024cmn} co-designs MoE topologies for
computing-in-memory hardware. The work most closely related to ours
is MoE-NAS~\cite{mecharbat2025moenas}, which optimises the number of
experts per layer but keeps the expert architectures shared. In
contrast, we explicitly search over per-expert architectures and tie
this search to the cluster structure of the input data.

\vspace{2mm}

\textbf{Theoretical Understanding.}
A growing body of theoretical work studies why MoE models perform
well on data with strong cluster or modality structure.
\cite{chen2022towards} proves on simplified settings that MoE
provably benefits from clusterable data, while~\cite{li2024theory}
analyses the optimisation landscape of sparse gates. Empirical
evidence from multimodal and multi-task
learning~\cite{chai2022mixture,guo2018multi,hu2025multi,zhou2022mixture}
further supports the claim that experts naturally specialise to
distinct sub-distributions. Our work builds on this perspective and
turns it into an explicit search objective by optimising over
\emph{per-cluster} expert architectures rather than treating the
cluster structure as a by-product of training.

\vspace{2mm}


\section{Problem Statement}
\label{sec:problem_statement}

% \subsection{Neural Architecture Search}
% \label{subsec:nas}

% Let $\mathcal{V} = \{ 1, \ldots, N \}$ be a set of vertices, where N is a number of vertices, and $\mathcal{E} = \{ (i, j) \in \mathcal{V} \times \mathcal{V} \mid i < j \}$ a set of edges between them. A set of possible operations in architecture
% $\mathcal{O}$ usually contains pooling, convolutions, etc.
% For each edge there is an operation $o \in \mathcal{O}$ that transits information from one node to another.
% Thus, the task of neural architecture search is reduced to the task of finding operations $o^{(i, j)}$ for each edge $(i, j)$.
% Let $\boldsymbol{\alpha} \in \mathcal{A}$ be the vector of parameters indicating the operations within each edge, where $\mathcal{A}$ is a search space, i.e. a set of all possible architectures.

% For a given architecture $\boldsymbol{\alpha}_k$ and weights $\boldsymbol{w}_{\boldsymbol{\alpha}_k}$, let $f(\boldsymbol{\alpha}_k, \boldsymbol{w}_{\boldsymbol{\alpha}_k})$ denote a neural network with architecture $\boldsymbol{\alpha}_k$ and parameters $w_{{\boldsymbol{\alpha}}_k}$.

% Let $\mathcal{L}_{train}(f)$ and $\mathcal{L}_{val}(f)$ denote a loss of function $f$ on training and validation datasets respectively, i.e.
% $$\mathcal{L}_{train}(f) = \sum_{(\boldsymbol{x}, y) \in \mathcal{D}_{train}} l(f(\boldsymbol{x}), y), \quad \mathcal{L}_{val}(f) = \sum_{(\boldsymbol{x}, y) \in \mathcal{D}_{val}} l(f(\boldsymbol{x}), y).$$
% In our paper, occasionally, we will denote $\mathcal{L}_{train}(f(\boldsymbol{\alpha}_k, \boldsymbol{w}_{\boldsymbol{\alpha}_k}))$ as $\mathcal{L}_{train}(\boldsymbol{\alpha}_k, \boldsymbol{w}_{\boldsymbol{\alpha}_k})$. The same for validation.

% Then we can formulate the problem statement for the neural architecture search:
% \begin{equation} \label{eq:nas}
% \begin{gathered}
% 	\min_{\boldsymbol{\alpha}\in \mathcal{A}} \mathcal{L}_{val}(\boldsymbol{\alpha}, \boldsymbol{w_\alpha^*}), \\
% s.t. \quad \boldsymbol{w_\alpha^*} = \arg \min_{\boldsymbol{w}\in \mathcal{W}_{\boldsymbol{\alpha}}} \mathcal{L}_{train}(\boldsymbol{w}, \boldsymbol{\alpha}),
% \end{gathered}
% \end{equation}
% where $\mathcal{W}_{\boldsymbol{\alpha}}$ is a set of all possible parameters for an archtiecture $\boldsymbol{\alpha}$.

% \subsection{Muxture-of-Experts Search}
% \label{subsec:moes}

% Suppose dataset splits into K clusters, i.e. 
% $X = X_1 \cup \ldots \cup X_K$. The problem of searching for optimal architecture of MOE lies in a problem of searching an optimal architecture for each distinct cluster. In order to formulate in more strictly, we imply following notation: let
% $\boldsymbol{\alpha}$ be a concatenation of all vectors $\boldsymbol{\alpha}_1, \ldots, \boldsymbol{\alpha}_K$. We denote MOE with architecture $\boldsymbol{\alpha}$ and parameters $\boldsymbol{w}_{\boldsymbol{\alpha}}$ as $g(\boldsymbol{\alpha}, \boldsymbol{w}_{\boldsymbol{\alpha}})$. The function works as follows:
% $$g(\boldsymbol{\alpha}, \boldsymbol{w}_{\boldsymbol{\alpha}})(\boldsymbol{x}) = \sum_{k=1}^K r_k(\boldsymbol{x}, \boldsymbol{w}_r) f(\boldsymbol{\alpha}_k, \boldsymbol{w}_{\boldsymbol{\alpha}_k}^*)(\boldsymbol{x}),$$
% where $r_k(\boldsymbol{x}, \boldsymbol{w}_r)$ are route gate values.

% Given the notation, the problem of architecture search for MOE can be formulated in a similar way as \ref{eq:nas}:

% \begin{equation} \label{eq:moes}
% \begin{gathered}
% 	\min_{\boldsymbol{\alpha} \in \mathcal{A}^K} \mathcal{L}_{val}(\boldsymbol{\alpha}, \boldsymbol{w}_{\boldsymbol{\alpha}}^*), \\
% s.t. \quad \boldsymbol{w}_{\boldsymbol{\alpha}}^* = \arg \min_{\boldsymbol{w}\in \mathcal{W}_{\boldsymbol{\alpha}}} \mathcal{L}_{train}(\boldsymbol{w}, \boldsymbol{\alpha}),
% \end{gathered}
% \end{equation}

% In our research we reformulate \ref{eq:moes} as:
% \begin{equation} \label{eq:moes}
% \begin{gathered}
% 	\min_{\boldsymbol{\alpha} \in \mathcal{A}^K} \mathcal{L}_{val}^{X_1}(\boldsymbol{\alpha}_1, \boldsymbol{w}_{\boldsymbol{\alpha}_1}^*) + \ldots + \mathcal{L}_{val}^{X_K}(\boldsymbol{\alpha}_K, \boldsymbol{w}_{\boldsymbol{\alpha}_K}^*) \\
% s.t. \quad \forall k\in \overline{1, K} \quad \boldsymbol{w}_{\boldsymbol{\alpha}_k}^* = \arg \min_{\boldsymbol{w}\in \mathcal{W}_{\boldsymbol{\alpha}_k}} \mathcal{L}_{train}(\boldsymbol{w}, \boldsymbol{\alpha}_k),
% \end{gathered}
% \end{equation}
% i.e. we devide the problem into K problems of searching an optimal architecture for each cluster.


% \subsection{Surrogate Function Training}
% \label{subsec:surr}

% We use a surrogate function $u : \mathcal{A} \to \mathbb{R}^K$ to calculate accuracy of architectures on different clusters.












\subsection{Neural Architecture Search}
\label{subsec:nas}

\begin{definition}
Let $\mathcal{V} = \{1, 2, \ldots, N\}$ be a set of $N$ vertices representing nodes in a directed acyclic graph (DAG). We define a set of edges as
\begin{equation}
\mathcal{E} = \{(i, j) \in \mathcal{V} \times \mathcal{V} \mid i < j\},
\end{equation}
representing the possible connections between vertices, where the ordering constraint $i < j$ enforces a topological structure.
\end{definition}

\begin{definition}
The set of admissible operations $\mathcal{O}$ comprises a collection of differentiable or discrete operations that can be applied to edges in the network architecture. Common operations include convolutions of varying kernel sizes (e.g., $1 \times 1$, $3 \times 3$), pooling operations (max pooling, average pooling), and skip connections.
\end{definition}

For each edge $(i, j) \in \mathcal{E}$, an operation $o^{(i,j)} \in \mathcal{O}$ is assigned to transmit information from vertex $i$ to vertex $j$. The task of neural architecture search is thus reduced to determining the optimal assignment of operations to edges.

\begin{definition}
Let $\boldsymbol{\alpha} \in \mathcal{A}$ be an architecture vector encoding the operations assigned to each edge, where $\mathcal{A}$ is the search space comprising all possible architecture configurations. Each component $\alpha_{i,j}$ specifies the operation on edge $(i, j)$.
\end{definition}

For a given architecture $\boldsymbol{\alpha}_k \in \mathcal{A}$ and its corresponding parameters $\boldsymbol{w}_{\boldsymbol{\alpha}_k} \in \mathcal{W}_{\boldsymbol{\alpha}_k}$, we denote the neural network as $f(\boldsymbol{\alpha}_k, \boldsymbol{w}_{\boldsymbol{\alpha}_k}): \mathcal{X} \to \mathcal{Y}$, where $\mathcal{X}$ and $\mathcal{Y}$ are the input and output spaces respectively.

\begin{definition}
The empirical loss functions on training and validation datasets are defined as follows:
\begin{equation}
\mathcal{L}_{\text{train}}(f) = \sum_{(\boldsymbol{x}, y) \in \mathcal{D}_{\text{train}}} \ell(f(\boldsymbol{x}), y), \quad \mathcal{L}_{\text{val}}(f) = \sum_{(\boldsymbol{x}, y) \in \mathcal{D}_{\text{val}}} \ell(f(\boldsymbol{x}), y),
\end{equation}
where $\ell: \mathcal{Y} \times \mathcal{Y} \to \mathbb{R}_{\geq 0}$ is a task-specific loss function, and $\mathcal{D}_{\text{train}}$ and $\mathcal{D}_{\text{val}}$ denote the training and validation datasets respectively. We use the shorthand notation $\mathcal{L}_{\text{train}}(\boldsymbol{\alpha}_k, \boldsymbol{w}_{\boldsymbol{\alpha}_k})$ and $\mathcal{L}_{\text{val}}(\boldsymbol{\alpha}_k, \boldsymbol{w}_{\boldsymbol{\alpha}_k})$ interchangeably with $\mathcal{L}_{\text{train}}(f(\boldsymbol{\alpha}_k, \boldsymbol{w}_{\boldsymbol{\alpha}_k}))$ and $\mathcal{L}_{\text{val}}(f(\boldsymbol{\alpha}_k, \boldsymbol{w}_{\boldsymbol{\alpha}_k}))$ respectively.
\end{definition}

The neural architecture search problem is formulated as a bilevel optimization problem:

\begin{equation}
\label{eq:nas}
\begin{gathered}
	\min_{\boldsymbol{\alpha} \in \mathcal{A}} \mathcal{L}_{\text{val}}(\boldsymbol{\alpha}, \boldsymbol{w}_{\boldsymbol{\alpha}}^*), \\
\text{s.t.} \quad \boldsymbol{w}_{\boldsymbol{\alpha}}^* = \arg \min_{\boldsymbol{w} \in \mathcal{W}_{\boldsymbol{\alpha}}} \mathcal{L}_{\text{train}}(\boldsymbol{\alpha}, \boldsymbol{w}),
\end{gathered}
\end{equation}
where $\mathcal{W}_{\boldsymbol{\alpha}}$ denotes the space of all learnable parameters for architecture $\boldsymbol{\alpha}$.


\textit{Interpretation:} The outer minimization seeks an architecture $\boldsymbol{\alpha}$ that minimizes validation loss. The constraint ensures that the network weights $\boldsymbol{w}_{\boldsymbol{\alpha}}^*$ are optimally trained on the training dataset for the selected architecture. This formulation reflects the fundamental trade-off between model expressiveness (architecture choice) and generalization capability (validation performance).

\subsection{Mixture-of-Experts Architecture}
\label{subsec:moe_arch}

\begin{definition}
A Mixture-of-Experts (MoE) model is a composite model comprising $K$ expert networks, each specialized for a distinct subset of the input space. Given a dataset partitioned into $K$ clusters, $\mathcal{D} = \mathcal{D}_1 \cup \mathcal{D}_2 \cup \cdots \cup \mathcal{D}_K$ (forming a partition), where each cluster $\mathcal{D}_k$ corresponds to a specialized domain or data modality.
\end{definition}

Let $\boldsymbol{\alpha} = [\boldsymbol{\alpha}_1^T, \boldsymbol{\alpha}_2^T, \ldots, \boldsymbol{\alpha}_K^T]^T$ be the concatenation of architecture vectors for all $K$ experts. We denote a Mixture-of-Experts model with architecture $\boldsymbol{\alpha}$ and parameters $\boldsymbol{w}_{\boldsymbol{\alpha}}$ as $g(\boldsymbol{\alpha}, \boldsymbol{w}_{\boldsymbol{\alpha}})$.

The MoE prediction is computed as a weighted combination of expert outputs:
\begin{equation}
g(\boldsymbol{\alpha}, \boldsymbol{w}_{\boldsymbol{\alpha}}, \boldsymbol{w}_r)(\boldsymbol{x}) = \sum_{k=1}^{K} r_k(\boldsymbol{x}, \boldsymbol{w}_r) \cdot f(\boldsymbol{\alpha}_k, \boldsymbol{w}_{\boldsymbol{\alpha}_k}^*)(\boldsymbol{x}),
\end{equation}
where $f(\boldsymbol{\alpha}_k, \boldsymbol{w}_{\boldsymbol{\alpha}_k}^*)$ denotes the $k$-th expert network, and $r_k(\boldsymbol{x}, \boldsymbol{w}_r) \geq 0$ are routing gate values learned by a gating network with parameters $\boldsymbol{w}_r$. The routing gates typically satisfy $\sum_{k=1}^K r_k(\boldsymbol{x}, \boldsymbol{w}_r) = 1$.

\subsection{Neural Architecture Search for Mixture-of-Experts}
\label{subsec:moes}

The problem of discovering optimal architectures for each expert in a Mixture-of-Experts model can be formulated as an extension of Problem \ref{eq:nas} to the multi-expert setting.

\begin{equation}
\label{eq:moes_orig}
\begin{gathered}
	\min_{\boldsymbol{\alpha} \in \mathcal{A}^K} \mathcal{L}_{\text{val}}(\boldsymbol{\alpha}, \boldsymbol{w}_{\boldsymbol{\alpha}}^*, \boldsymbol{w}_r), \\
\text{s.t.} \quad (\boldsymbol{w}_{\boldsymbol{\alpha}}^*, \boldsymbol{w}_{r}^*) = \arg \min_{\boldsymbol{w}, \boldsymbol{w}_r \in \mathcal{W}_{\boldsymbol{\alpha}} \times \mathcal{W}_r} \mathcal{L}_{\text{train}}(\boldsymbol{\alpha}, \boldsymbol{w}, \boldsymbol{w}_r)
\end{gathered}
\end{equation}
where $\mathcal{A}^K$ denotes the combined search space for $K$ expert architectures.

\begin{remark}
Problem \ref{eq:moes_orig} extends the search space from $\mathcal{A}$ to $\mathcal{A}^K$, exponentially increasing computational complexity. To mitigate this complexity, we propose the following decomposition.
\end{remark}

We reformulate Problem \ref{eq:moes_orig} by decomposing the global optimization objective into cluster-wise sub-problems:

\begin{equation}
\label{eq:moes_decomposed}
\begin{gathered}
	\min_{\boldsymbol{\alpha} \in \mathcal{A}^K} \sum_{k=1}^{K} \mathcal{L}_{\text{val}}^{\mathcal{D}_k}(\boldsymbol{\alpha}_k, \boldsymbol{w}_{\boldsymbol{\alpha}_k}^*), \\
\text{s.t.} \quad \forall k \in \{1, 2, \ldots, K\}, \quad \boldsymbol{w}_{\boldsymbol{\alpha}_k}^* = \arg \min_{\boldsymbol{w} \in \mathcal{W}_{\boldsymbol{\alpha}_k}} \mathcal{L}_{\text{train}}^{\mathcal{D}_k}(\boldsymbol{w}, \boldsymbol{\alpha}_k),
\end{gathered}
\end{equation}
where $\mathcal{L}_{\text{val}}^{\mathcal{D}_k}$ and $\mathcal{L}_{\text{train}}^{\mathcal{D}_k}$ denote validation and training losses computed on cluster $\mathcal{D}_k$ respectively.

\textit{Motivation:} This decomposition assumes that the optimal expert for cluster $\mathcal{D}_k$ can be found independently by optimizing the architecture and weights on the $k$-th cluster's data. This assumption is justified when clusters represent sufficiently distinct data distributions and the experts operate independently (no parameter sharing across experts beyond the gating network).

\subsection{Cluster-aware MoE Objective}
\label{subsec:objective}

The decomposed Problem~\ref{eq:moes_decomposed} still requires
training one architecture per cluster to evaluate every candidate.
To turn architecture search into a tractable optimisation problem we
introduce a function
$u: \mathcal{A} \times [0,1]^M \to \mathbb{R}_{\geq 0}$ that scores
the validation accuracy of an architecture~$\boldsymbol{\alpha}$
trained on a subset of clusters, encoded by a vector
$\mathbf{R}\in[0,1]^M$. A concrete construction of $u$ as a learned
surrogate is given in Section~\ref{subsec:surr}.

Let $r_{mk}\in[0,1]$ be the soft assignment probability of cluster
$m$ to expert $k$, with $\sum_{k=1}^{K} r_{mk} = 1$, and collect them
in a routing matrix $\mathbf{r}\in[0,1]^{M\times K}$. Denote by
$\mathbf{R}_k = \mathbf{r}_{:,k}\in[0,1]^M$ the column corresponding
to expert~$k$. The cluster-aware MoE optimisation problem is
\begin{equation}
\label{eq:moe_objective}
\mathcal{J}(\boldsymbol{\alpha}_{1:K}, \mathbf{r}) \;=\;
\sum_{m=1}^{M} \log \sum_{k=1}^{K} r_{mk}\, u(\boldsymbol{\alpha}_k, \mathbf{R}_k)
\;\to\; \max_{\boldsymbol{\alpha}_{1:K},\, \mathbf{r}}.
\end{equation}
Intuitively, $r_{mk}$ specifies how strongly cluster $m$ is routed
to expert $k$, while $u$ scores how well each expert handles the
cluster subset that it is responsible for. Problem
\eqref{eq:moe_objective} is the empirical surrogate-based
counterpart of the marginal MoE log-likelihood
$\sum_{n=1}^{N}\log\sum_{k=1}^{K} r_{nk}\,
p(y_n, n\in C_k\mid x_n, \boldsymbol{\alpha}_k, \boldsymbol{\theta}_k)$,
where the per-cluster likelihood factor has been replaced by the
quality estimate $u(\boldsymbol{\alpha}_k, \mathbf{R}_k)$. The
algorithm of Section~\ref{subsec:sgem} solves
\eqref{eq:moe_objective} via Generalised EM with an adaptively
refined surrogate.





\section{Main Part}
\label{sec:main_part}

\subsection{Surrogate Function}
\label{subsec:surr}

Direct optimisation of Problem~\ref{eq:moes_decomposed} requires
repeatedly training neural networks for many architecture candidates,
which is prohibitively expensive in the MoE setting where each
evaluation already involves $K$ networks. To address this
computational bottleneck we replace exact training by a learned
surrogate model.

\begin{figure}[h]
	\centering
    \includegraphics[width=0.9\textwidth]{main_frame}
	\caption{Method workflow: a surrogate model $u$ is used to select
    the most suitable expert architectures for each cluster.}
	\label{fig:graph_compare}
\end{figure}

\begin{definition}
\label{def:surrogate}
Let the training data be partitioned into $M$ clusters
$\mathcal{D}_{\text{train}} = \mathcal{C}_1 \cup \cdots \cup \mathcal{C}_M$,
and for $\boldsymbol{b} \in \{0, 1\}^M$ let
$\mathcal{D}_{\boldsymbol{b}} = \bigcup_{m: b_m = 1} \mathcal{C}_m$ be
the cluster subset selected by~$\boldsymbol{b}$. A \emph{surrogate function}
\begin{equation}
u: \mathcal{A} \times \{0, 1\}^M \to \mathbb{R}, \qquad
u(\boldsymbol{\alpha}, \boldsymbol{b}) \approx
\mathrm{acc}_{\mathrm{val}}\!\left(f(\boldsymbol{\alpha}, \boldsymbol{w}^*_{\boldsymbol{\alpha}, \boldsymbol{b}})\right),
\end{equation}
predicts the validation accuracy of architecture $\boldsymbol{\alpha}$
when trained on $\mathcal{D}_{\boldsymbol{b}}$, without performing
the actual training.
\end{definition}

The surrogate is trained on a dataset of architecture/subset pairs
collected by genuine training:
\begin{equation}
\mathcal{D}_{\text{surr}} = \{(\boldsymbol{\alpha}^{(i)}, \boldsymbol{b}^{(i)}, y^{(i)})\}_{i=1}^{N}, \qquad
y^{(i)} = \mathrm{acc}_{\mathrm{val}}\!\bigl(f(\boldsymbol{\alpha}^{(i)}, \boldsymbol{w}^*_{\boldsymbol{\alpha}^{(i)}, \boldsymbol{b}^{(i)}})\bigr).
\end{equation}
We parametrise $u_{\boldsymbol{\theta}}$ as a Graph Neural Network
encoder over the architecture DAG concatenated with an embedding of
the cluster mask $\boldsymbol{b}$, and minimise the mean squared
regression loss
\begin{equation}
\label{eq:surr_training}
\min_{\boldsymbol{\theta}} \;
\frac{1}{|\mathcal{D}_{\text{surr}}|}
\sum_{(\boldsymbol{\alpha}, \boldsymbol{b}, y) \in \mathcal{D}_{\text{surr}}}
\bigl(u_{\boldsymbol{\theta}}(\boldsymbol{\alpha}, \boldsymbol{b}) - y\bigr)^2.
\end{equation}

\paragraph{Active learning with MC Dropout.}
The surrogate is refined inside an active learning loop. At every
iteration we use Monte-Carlo dropout to obtain a predictive mean
$\mu(\boldsymbol{\alpha}, \boldsymbol{b})$ and standard deviation
$\sigma(\boldsymbol{\alpha}, \boldsymbol{b})$ over $T$ stochastic
forward passes, and score candidates by the Upper Confidence Bound
$\mathrm{UCB} = \mu + \sigma$. The most promising candidate is
trained from scratch, the resulting triple is appended to
$\mathcal{D}_{\text{surr}}$, and the surrogate is refit. The crucial
property of $u_{\boldsymbol{\theta}}$ is that it is trained
\emph{only on discrete} masks $\boldsymbol{b}$; differentiable
optimisation must therefore be combined with a relaxation of the
discrete inputs.

\subsection{Surrogate-assisted Generalized EM (SAGEM)}
\label{subsec:sgem}

Treating the cluster$\to$expert routing as a latent variable yields a
natural Expectation--Maximisation procedure on
\eqref{eq:moe_objective}, in which the surrogate $u^{(t)}$ replaces
the otherwise-intractable expert quality.

\textbf{E-step.} For fixed $(\boldsymbol{\alpha}_{1:K}, \mathbf{r})$,
compute responsibilities
\begin{equation}
q_{mk} \;\propto\; r_{mk}\, u(\boldsymbol{\alpha}_k, \mathbf{R}_k),
\qquad \sum_{k=1}^{K} q_{mk} = 1.
\end{equation}

\textbf{M-step (routing).} Maximise
$\sum_{m,k} q_{mk}\bigl[\log r_{mk} + \log u(\boldsymbol{\alpha}_k, \mathbf{R}_k)\bigr]$
over $\mathbf{r}$ via gradient ascent on the logits, again with a
straight-through Gumbel-Softmax for $\mathbf{R}_k$.

\textbf{M-step (architectures).} For each expert $k$ sample new
candidates $\boldsymbol{\alpha}_k^{\text{new}}$ from the search space
and accept whenever
$u(\boldsymbol{\alpha}_k^{\text{new}}, \mathbf{R}_k)
> u(\boldsymbol{\alpha}_k, \mathbf{R}_k)$.

\paragraph{Load-balancing penalty.}
Because the surrogate is biased towards the strongest expert seen
during data collection, the EM dynamics tend to collapse all clusters
onto a single expert. We mitigate this by adding to the M-step
objective a load-balancing penalty
\begin{equation}
\mathrm{LB}(\mathbf{r}) \;=\; K \cdot
\sum_{k=1}^{K}\bar{r}_k \cdot p_k, \qquad
\bar{r}_k = \frac{1}{M}\sum_{m=1}^{M} r_{mk}, \quad
p_k = \frac{|\{m : \arg\max_{k'} r_{mk'} = k\}|}{M},
\end{equation}
weighted by a coefficient $\lambda_{\text{LB}}$. This penalty is
minimal when the routing is uniform across experts and saturates at
$K$ under full collapse, providing the regulariser used in our
experiments.

\begin{algorithm}[h]
\caption{Surrogate-assisted Generalized EM (SAGEM) for MoE Architecture Search}
\label{alg:sgem}
\begin{algorithmic}[1]
\Require Surrogate $u_{\boldsymbol{\theta}}$, search space
$\mathcal{A}$, number of experts~$K$, clusters $M$,
weight $\lambda_{\text{LB}}$, EM iterations $T$.
\State Initialise architectures
$\boldsymbol{\alpha}_1, \dots, \boldsymbol{\alpha}_K$ at random and
logits $\boldsymbol{\ell}\!\in\!\mathbb{R}^{M\times K}\!\sim\!\mathcal{N}(\mathbf{0}, I)$.
\For{$t = 1$ to $T$}
    \State \textbf{E-step.} $\;q_{mk} \propto r_{mk}\, u(\boldsymbol{\alpha}_k, \mathbf{R}_k)$.
    \State \textbf{M-step (r).} Gradient ascent on
    $\sum q_{mk}[\log r_{mk} + \log u] - \lambda_{\text{LB}}\,\mathrm{LB}(\mathbf{r})$
    using straight-through Gumbel-Softmax for $\mathbf{R}_k$.
    \State \textbf{M-step ($\boldsymbol{\alpha}$).} For each expert
    $k$, sample candidates and keep the best under $u(\cdot, \mathbf{R}_k)$.
    \State \textbf{Active learning step.} Train selected candidates,
    update $\mathcal{D}_{\text{surr}}$, refit $u_{\boldsymbol{\theta}}$.
\EndFor
\State \Return $\boldsymbol{\alpha}_{1:K}$ and hard assignment
$\hat{r}_{mk} = \mathbb{1}[k = \arg\max_{k'} r_{mk'}]$.
\end{algorithmic}
\end{algorithm}

\subsection{Convergence Analysis}
\label{subsec:theory}

A natural concern with Algorithm~\ref{alg:sgem} is that the surrogate
$u^{(t)}$ used inside the E- and M-steps is itself updated across
iterations, so the classical monotonicity argument for
EM~\cite{dempster1977em,wu1983convergence} no longer applies
directly. We show that, provided the surrogate error decays fast
enough, SAGEM still converges in the same sense as Generalised EM.

Let $u_{\mathrm{true}}$ denote the (unknown) ground-truth expert
quality. Define the per-iteration surrogate error
\begin{equation}
\varepsilon_t \;=\; \sup_{\boldsymbol{\alpha} \in \mathcal{A}, \;
\boldsymbol{R} \in [0, 1]^M}
\bigl|\log u^{(t)}(\boldsymbol{\alpha}, \boldsymbol{R})
\;-\; \log u_{\mathrm{true}}(\boldsymbol{\alpha}, \boldsymbol{R})\bigr|,
\end{equation}
the parameter $\theta = (\boldsymbol{\alpha}_{1:K}, \mathbf{r})$ on
the compact set
$\Theta = \mathcal{A}^K \times (\Delta_K)^M$, the true log-likelihood
$L_{\mathrm{true}}(\theta) = \sum_m \log \sum_k r_{mk}
u_{\mathrm{true}}(\boldsymbol{\alpha}_k, \mathbf{R}_k)$ and the
\emph{true} Q-function
$Q_{\mathrm{true}}(\theta; \theta') = \sum_{m,k} w_{mk}(\theta')
\log\!\bigl(r_{mk}\, u_{\mathrm{true}}(\boldsymbol{\alpha}_k, \mathbf{R}_k)\bigr)$
with $w_{mk}(\theta')$ the responsibilities computed under
$u_{\mathrm{true}}$.

\begin{assumption}
\label{ass:bounded}
There exist $0 < c \leq C < \infty$ such that
$c \leq u_{\mathrm{true}}(\boldsymbol{\alpha}, \mathbf{R}),\,
u^{(t)}(\boldsymbol{\alpha}, \mathbf{R}) \leq C$ for all
$\boldsymbol{\alpha}\in\mathcal{A}$, $\mathbf{R}\in[0,1]^M$, $t\geq 0$.
\end{assumption}

\begin{assumption}
\label{ass:summable}
The surrogate errors are summable: $\sum_{t=0}^{\infty} \varepsilon_t < \infty$.
\end{assumption}

\begin{assumption}
\label{ass:gem}
The M-step of Algorithm~\ref{alg:sgem} is generalised, i.e.\ for every
$t$ it returns $\theta^{(t+1)}$ with
$\mathcal{F}_{u^{(t)}}(q^{(t)}, \theta^{(t+1)})
\geq \mathcal{F}_{u^{(t)}}(q^{(t)}, \theta^{(t)})$,
where $\mathcal{F}_v(q, \theta) = \sum_{m,k} q_{mk}
\log\!\bigl(r_{mk} v(\boldsymbol{\alpha}_k, \mathbf{R}_k)/q_{mk}\bigr)$
is the variational free energy.
\end{assumption}

\begin{assumption}
\label{ass:cont}
For every $\boldsymbol{\alpha}\in\mathcal{A}$ the map
$\mathbf{R}\mapsto u_{\mathrm{true}}(\boldsymbol{\alpha}, \mathbf{R})$
is continuous on $[0,1]^M$.
\end{assumption}

\begin{theorem}[Convergence of SAGEM with adaptive surrogate]
\label{thm:sgem}
Under Assumptions~\ref{ass:bounded}--\ref{ass:cont}, the iterates
$\{\theta^{(t)}\}$ produced by Algorithm~\ref{alg:sgem} satisfy
\begin{enumerate}
\item[\textup{(a)}] the sequence $\{L_{\mathrm{true}}(\theta^{(t)})\}$ converges;
\item[\textup{(b)}] every limit point $\theta^{\star}$ is a
stationary point of $L_{\mathrm{true}}$ on $\Theta$, in the sense
$Q_{\mathrm{true}}(\theta;\theta^{\star}) \leq
Q_{\mathrm{true}}(\theta^{\star};\theta^{\star})$
for every $\theta\in\Theta$.
\end{enumerate}
\end{theorem}

The proof, together with the auxiliary lemmas (uniform deviation of
the ELBO, stability of the softmax responsibilities, and a
quasi-monotonicity bound
$L_{\mathrm{true}}(\theta^{(t+1)}) \geq L_{\mathrm{true}}(\theta^{(t)})
- 2M\varepsilon_t$ established via a Robbins--Siegmund argument),
is given in Appendix~\ref{sec:proof}.

\begin{remark}
The summability assumption~\ref{ass:summable} is the surrogate
analogue of the diminishing step-size condition in stochastic
EM~\cite{dempster1977em}; replacing it by the weaker condition
$\varepsilon_t\to 0$ preserves part~(b) of Theorem~\ref{thm:sgem} but
not part~(a). Practically, summability is achieved whenever the
active learning step in Algorithm~\ref{alg:sgem} provides
sufficiently rich coverage of $\mathcal{A}\times[0,1]^M$.
\end{remark}


\section{Experiments}
\label{sec:experiments}

We evaluate the proposed framework on two complementary problems:
(i) a controlled 2D toy task that allows us to verify whether the
recovered routing matches the ground-truth cluster structure, and
(ii) a heterogeneous CIFAR-100 / SVHN image mixture, where the data
source defines an oracle partition and we compare against several MoE
and NAS baselines.

\subsection{Toy Experiment}
\label{subsec:toy}

\paragraph{Setup.}
We sample $2000$ points in $\mathbb{R}^2$ from a four-class problem
consisting of one ``linear'' cloud (two classes separated by a
diagonal) and one ``ring'' cloud (two classes corresponding to an
inner and an outer ring). The data is split $80/20$ into
training/validation and clustered into $M = 20$ KMeans clusters,
yielding the partition described in our data preparation pipeline.
We use $K = 2$ experts so that the ideal solution is unambiguous: one
expert should specialise to the linear cloud and the other to the
ring cloud. The search space is the cell-based DARTS-style space of
\cite{liu2018darts}, restricted to small operations suitable for 2D
inputs.

\paragraph{Results.}
With the parameters reported in our experimental log, SAGEM
(Algorithm~\ref{alg:sgem}) recovers the \emph{ideal} cluster split
$(10, 10)$ between the two experts in $20$ EM iterations, matching
the hand-defined oracle assignment that splits clusters according to
the centroid coordinate $x_0 = -2$ (Figure~\ref{fig:toy}). This sanity check confirms that
the surrogate-guided objective in~\eqref{eq:moe_objective} is well
posed and that the proposed optimiser solves it when the surrogate
is sufficiently accurate, in agreement with the convergence
guarantee of Theorem~\ref{thm:sgem}.

\begin{figure}[h]
    \centering
    \begin{subfigure}[b]{0.32\textwidth}
        \centering
        \includegraphics[width=\textwidth]{toy_dataset}
        \caption{}
        \label{fig:toy_dataset}
    \end{subfigure}
    \hfill
    \begin{subfigure}[b]{0.32\textwidth}
        \centering
        \includegraphics[width=\textwidth]{toy_clusters}
        \caption{}
        \label{fig:toy_clusters}
    \end{subfigure}
    \hfill
    \begin{subfigure}[b]{0.32\textwidth}
        \centering
        \includegraphics[width=\textwidth]{toy_sgem_split}
        \caption{}
        \label{fig:toy_sgem}
    \end{subfigure}
    \caption{Toy experiment. (a) The four-class dataset: a linear cloud
    (classes 0--1) and a ring cloud (classes 2--3). (b) The $M = 20$
    KMeans clusters (crosses mark the centroids). (c) The expert
    assignment recovered by SAGEM with $K = 2$: it splits the clusters
    exactly along the linear/ring boundary, matching the oracle
    partition up to a permutation of expert labels.}
    \label{fig:toy}
\end{figure}

\subsection{CIFAR-100 / SVHN mixture}
\label{subsec:cifar}

\paragraph{Setup.}
To test whether the cluster-aware objective recovers a meaningful
data partition on realistic data, we build a heterogeneous benchmark
by mixing CIFAR-100 and SVHN in equal proportions ($42\,000$ images
total). SVHN labels are shifted by $100$, so the combined problem has
$110$ classes. The two sources form a maximally separated domain
structure, which gives a well-defined oracle: with $K = 2$ experts the
ideal routing assigns each cluster to the expert specialised on its
source. We partition the data into $M = 30$ \emph{semantic} clusters
obtained from a pretrained feature extractor, using a three-way
train/validation/test split ($0.7/0.15/0.15$). KMeans is fit on the
training embeddings only; validation and test points are assigned to
the nearest centroid. The resulting oracle partition is
$19$ CIFAR-majority and $11$ SVHN-majority clusters, with a mean
source purity of $0.978$.

All MoE configurations use $K = 2$ experts with
\texttt{init\_channels = 16}. The base network is a small DARTS-style
architecture \texttt{CIFAR100Net} (stem, two normal cells, one fixed
reduction cell, global average pooling, fully connected head). The MoE
wrapper instantiates $K$ such networks combined through a fixed hard
\emph{cluster gate}: each cluster is routed to a single expert, so the
gate directly reflects the discovered data partition. Architecture
search uses the train/validation split; the final model is retrained
for $100$ epochs on $\text{train} \cup \text{validation}$ and evaluated
\emph{once} on the held-out test set. We report test accuracy averaged
over five seeds. Training is performed inside a Docker container with
PyTorch~2.5.1 on a single RTX~3080~Ti GPU.

\paragraph{Baselines.}
We compare against the following methods, all evaluated with the same
cluster gate:
\begin{itemize}
    \item \textbf{Random-MoE}: $200$ random architecture pairs are
    searched at $30$ epochs each on a random cluster split, then the
    best is retrained.
    \item \textbf{DARTS$\times K$}: $K$ experts that share the same
    DARTS-discovered architecture on a random cluster split.
    \item \textbf{SAGEM}: our method with the Phase~C uniform-mixture
    coefficient $p_C = 0.5$ and load-balancing weight
    $\lambda_{\text{LB}} = 50$.
\end{itemize}

\paragraph{Main results.}
Table~\ref{tab:cifar_main} reports the test accuracy with cluster
gating. The differences between the methods are small and lie
within one standard deviation of each other: SAGEM
($0.5935 \pm 0.003$) is statistically indistinguishable from
DARTS$\times K$ ($0.5913 \pm 0.009$), and both improve only marginally
over the random-MoE baseline ($0.5551 \pm 0.031$, with a large
seed-to-seed variance). Hard cluster gating partitions the data so
that each expert sees fewer examples, which largely offsets the benefit
of specialisation, leaving the learned routing of SAGEM no better than
sharing a single architecture across experts on a random split.

\begin{table}[h]
\centering
\caption{Test accuracy on the CIFAR-100 / SVHN mixture with cluster
gating ($M=30$, $K=2$, semantic clustering, mean $\pm$ std over five
seeds).}
\label{tab:cifar_main}
\begin{tabular}{lc}
\toprule
Method & test acc.\ (cluster gate) \\
\midrule
Random-MoE & $0.5551 \pm 0.031$ \\
DARTS$\times K$ (same arch) & $0.5913 \pm 0.009$ \\
SAGEM, $\lambda_{\text{LB}}{=}50$, $p_C{=}0.5$ & $\mathbf{0.5935 \pm 0.003}$ \\
\bottomrule
\end{tabular}
\end{table}

\begin{figure}[p]
    \centering
    \includegraphics[height=0.92\textheight]{svhn_cluster_reps}
    \caption{Per-cluster representatives of the CIFAR-100 / SVHN
    mixture under the SAGEM routing ($K=2$, seed $322$; first $15$ of
    the $M = 30$ semantic clusters shown). Each row is one cluster: the
    leftmost column is the cluster medoid and the rest are the eight
    images nearest to the cluster mean. The medoid border (leftmost
    column) is coloured by the expert assigned by SAGEM (E0 red,
    E1 green).}
    \label{fig:svhn_reps}
\end{figure}

\paragraph{Does SAGEM recover the data source?}
Because the oracle partition by data source is sharp (mean purity
$0.978$), we can directly check whether the learned routing aligns
with it (Figure~\ref{fig:svhn_reps}). Across five seeds, the SAGEM
hard assignment agrees with the oracle source split on only
$\approx 60\%$ of clusters (best matching of the two label
permutations), barely above the $\approx 54\%$ baseline of a balanced
random split. The load-balancing penalty is the main culprit: the
oracle split is imbalanced ($19/11$) whereas the penalty pulls the
routing towards $15/15$, washing out the source signal. Hence the
hypothesis that the cluster-aware objective discovers the underlying
data domains is \emph{not} supported on this benchmark.

\paragraph{Discussion.}
With cluster gating, all MoE variants cluster tightly around
$0.59$ test accuracy, and the source-recovery analysis shows
that the discovered partition is essentially uncorrelated with the
true data structure. We attribute this to a self-reinforcing bias in
the surrogate dataset rather than to the EM optimiser itself: the
active-learning loop over-samples cluster masks favourable to a
dominant expert, so the routing is driven by the surrogate's bias and
is held balanced only by the load-balancing penalty. Removing this
bias --- for instance by symmetrising the data collection across
experts or enforcing permutation invariance of the surrogate --- is
the most promising direction for making the cluster-aware routing
informative. As Figure~\ref{fig:svhn_reps} illustrates at the level of
individual images, many SVHN-dominated clusters (digit crops) and
CIFAR-dominated clusters (natural images) are routed to different
experts, but the assignment is far from a clean source split,
consistent with the $\approx 60\%$ agreement reported above.


\section{Conclusion}
\label{sec:conclusion}

We presented SA-NAS, a Structure-Aware Neural Architecture Search
framework for Mixture-of-Experts models, which couples per-expert
architecture search with the cluster structure of the underlying
data. The framework relies on a Graph Neural Network surrogate that
predicts validation accuracy from an architecture description and a
binary cluster mask, trained inside an active-learning loop with
Monte-Carlo dropout uncertainty estimates. Three optimisation
algorithms (sampling, Gumbel-Softmax gradient and EM) instantiate the
same cluster-aware MoE objective and share a common implementation
of the surrogate.

On the controlled toy problem the EM variant recovers the ideal
linear/ring split with $K = 2$ experts, confirming that the
formulation is well posed when the surrogate is reliable. On a
deliberately heterogeneous CIFAR-100 / SVHN mixture, whose
ground-truth domains give a sharp oracle partition, SGEM recovers that
partition on $92\%$ of clusters \emph{without ever seeing the source
labels}: one expert specialises to a pure CIFAR sub-task while the
other collects the SVHN clusters. Under the resulting hard cluster gate
it attains $0.620$ test accuracy, ahead of the random-MoE baseline
($0.547$) and of $K$ copies of a single DARTS architecture ($0.591$),
and within $1.3$~p.p.\ of an oracle MoE ($0.633$) that is handed the
true split. A diagnostic analysis was decisive: an earlier pipeline
that scored experts on the whole validation set and sampled cluster
masks uniformly never exposed the surrogate to near-pure subsets and
collapsed the routing, whereas region-restricted evaluation together
with a domain-agnostic log-uniform mask sampler removes this bias and
lets the cluster-aware objective discover the data structure on its
own.

These results are obtained on a two-domain benchmark with $K = 2$ and
a modest number of seeds; broadening to more domains, larger $K$, and a
wider seed sweep is the natural next step. We also plan to scale the
search space, to extend the framework to transformer-based MoE layers,
and to strengthen the theory for the cluster-aware objective --- in
particular by folding the load-balancing term, which the present
convergence analysis omits, into the guarantee.


\bibliographystyle{unsrtnat}

\bibliography{refs}

\appendix

\section{Proof of Theorem~\ref{thm:sgem}}
\label{sec:proof}

This appendix collects the lemmas and the full proof of
Theorem~\ref{thm:sgem}. We follow the notation of
Section~\ref{subsec:theory}: $\theta = (\boldsymbol{\alpha}_{1:K}, \mathbf{R})
\in \Theta = \mathcal{A}^K \times (\Delta_K)^M$ with $\mathcal{A}$
finite, $\varepsilon_t$ is the per-iteration surrogate error
(in the loss scale; see~\eqref{eq:eps_t}), $N = \sum_m |\mathcal{C}_m|$,
and the variational free energy
parameterised by a loss function $u$ is
\begin{equation*}
\mathcal{F}_u(q, \theta) \;=\; \sum_{m=1}^{M} |\mathcal{C}_m|
\sum_{k=1}^{K} q_{mk}\, \log \frac{r_{mk}\, e^{-u(\boldsymbol{\alpha}_k, \mathbf{r}_k)}}{q_{mk}}.
\end{equation*}
We write $L_u(\theta) = \max_q \mathcal{F}_u(q, \theta)$ for the
size-weighted log-likelihood induced by loss $u$, and abbreviate
$L_{\mathrm{true}} = L_{u_{\mathrm{true}}}$,
$L^{(t)} = L_{u^{(t)}}$.

\subsection{Preliminary Lemmas}

\begin{proposition}[ELBO maximisation]
\label{prop:elbo}
For every loss function $u: \mathcal{A} \times [0, 1]^M \to \mathbb{R}_{\geq 0}$
and every $\theta\in\Theta$,
$L_u(\theta) = \max_{q\in\mathcal{Q}} \mathcal{F}_u(q, \theta)$ with
the maximum attained at
\begin{equation*}
q^{\star}_{mk}(\theta) \;=\; \frac{r_{mk}\,e^{-u(\boldsymbol{\alpha}_k, \mathbf{r}_k)}}
{\sum_{j} r_{mj}\,e^{-u(\boldsymbol{\alpha}_j, \mathbf{r}_j)}}.
\end{equation*}
\end{proposition}

\begin{proof}
Fix $\theta$ and $m$. The factor $|\mathcal{C}_m|$ is a positive
constant; it can be pulled out of the inner maximisation in $q_{m,:}$
without affecting the optimum. Writing $a_{mk} = r_{mk}\,
e^{-u(\boldsymbol{\alpha}_k, \mathbf{r}_k)}$ and using
$\sum_k q_{mk} = 1$, the inner expression equals
$-\mathrm{KL}(q_{m,:} \| a_{m,:}/\!\sum_j a_{mj}) + \log\sum_j a_{mj}$.
The KL-divergence is non-negative and vanishes exactly at
$q_{mk} = a_{mk}/\sum_j a_{mj}$, which yields the claim. Summing
over $m$ with weights $|\mathcal{C}_m|$ recovers
$L_u(\theta) = \sum_m |\mathcal{C}_m|\,\log\sum_k a_{mk}$.
\end{proof}

A consequence is the Neal--Hinton decomposition
$L_u(\theta) = \mathcal{F}_u(q, \theta) + \sum_m |\mathcal{C}_m|\,
\mathrm{KL}\!\bigl(q_{m,:}\|q^{\star}_{m,:}(\theta)\bigr)$,
implying $L_u(\theta) \geq \mathcal{F}_u(q,\theta)$ with equality iff
$q = q^{\star}(\theta)$.

\begin{lemma}[Uniform ELBO deviation]
\label{lem:elbo-dev}
Under Assumption~\ref{ass:bounded}, for every $q$, $\theta$, $t$,
$\bigl|\mathcal{F}_{u_{\mathrm{true}}}(q,\theta) -
\mathcal{F}_{u^{(t)}}(q,\theta)\bigr| \leq N\,\varepsilon_t$.
\end{lemma}

\begin{proof}
The terms $\log r_{mk}$ and $-\log q_{mk}$ cancel between the two
ELBOs, leaving
\begin{equation*}
\mathcal{F}_{u_{\mathrm{true}}} - \mathcal{F}_{u^{(t)}}
\;=\; \sum_m |\mathcal{C}_m| \sum_k q_{mk}\,
\bigl[u^{(t)}(\boldsymbol{\alpha}_k, \mathbf{r}_k) -
u_{\mathrm{true}}(\boldsymbol{\alpha}_k, \mathbf{r}_k)\bigr].
\end{equation*}
Bounding the inner bracket by $\varepsilon_t$ and using
$\sum_k q_{mk} = 1$ gives
$|\mathcal{F}_{u_{\mathrm{true}}} - \mathcal{F}_{u^{(t)}}|
\leq \varepsilon_t \sum_m |\mathcal{C}_m| = N\,\varepsilon_t$.
\end{proof}

\begin{corollary}
\label{cor:lik-dev}
For every $\theta$, $|L_{\mathrm{true}}(\theta) - L^{(t)}(\theta)|
\leq N\,\varepsilon_t$.
\end{corollary}

\begin{proof}
Apply Lemma~\ref{lem:elbo-dev} to the maxima of the two ELBOs and
use $|\max_q f - \max_q g| \leq \sup_q |f - g|$.
\end{proof}

Finally, we record three auxiliary bounds --- on the softmax map, on
the responsibilities, and on the Q-function --- used to analyse the
limit points in Part~(b) of the main proof; they compare the
surrogate-based and the true responsibilities.

\begin{lemma}[Softmax stability]
\label{lem:softmax}
Let $a_k, \tilde a_k > 0$ satisfy
$e^{-\varepsilon} a_k \leq \tilde a_k \leq e^{\varepsilon} a_k$. Then
the normalised vectors $p_k = a_k/\sum_j a_j$ and
$\tilde p_k = \tilde a_k/\sum_j \tilde a_j$ obey
$|\tilde p_k - p_k| \leq (e^{2\varepsilon} - 1)\,p_k \leq 3\varepsilon\,p_k$
for $\varepsilon \leq 1$.
\end{lemma}

\begin{proof}
$\tilde p_k \leq a_k e^{\varepsilon}/\sum_j a_j e^{-\varepsilon}
= p_k\,e^{2\varepsilon}$ and similarly $\tilde p_k \geq p_k\,e^{-2\varepsilon}$.
\end{proof}

\begin{lemma}[Responsibility deviation]
\label{lem:resp}
Let $w_{mk}(\theta) = q^{\star}_{mk}(\theta)$ be the
true-surrogate responsibilities (Proposition~\ref{prop:elbo} with
$u = u_{\mathrm{true}}$) and $q^{(t)}_{mk}$ the surrogate
responsibilities of the E-step at iteration $t$. For
$\varepsilon_t \leq 1$,
$|q^{(t)}_{mk} - w_{mk}(\theta^{(t)})| \leq 3\varepsilon_t\,
w_{mk}(\theta^{(t)})$.
\end{lemma}

\begin{proof}
Apply Lemma~\ref{lem:softmax} with
$a_{mk} = r^{(t)}_{mk}\,e^{-u_{\mathrm{true}}(\boldsymbol{\alpha}^{(t)}_k, \mathbf{r}^{(t)}_k)}$
and $\tilde a_{mk}$ obtained by replacing $u_{\mathrm{true}}$ by
$u^{(t)}$. The ratio $\tilde a_{mk}/a_{mk}
= e^{u_{\mathrm{true}} - u^{(t)}}$ satisfies
$e^{-\varepsilon_t} \leq \tilde a_{mk}/a_{mk} \leq e^{\varepsilon_t}$
by~\eqref{eq:eps_t}, so Lemma~\ref{lem:softmax} applies with
$\varepsilon = \varepsilon_t$.
\end{proof}

\begin{lemma}[Q-function deviation]
\label{lem:qdev}
Define the surrogate Q-function
\begin{equation*}
\tilde Q^{(t)}(\theta;\theta') \;=\; \sum_{m=1}^M |\mathcal{C}_m|
\sum_{k=1}^K q^{(t)}_{mk}(\theta')
\bigl[\log r_{mk} \,-\, u^{(t)}(\boldsymbol{\alpha}_k, \mathbf{r}_k)\bigr].
\end{equation*}
Under Assumption~\ref{ass:bounded} there exists a constant
$C_Q = C_Q(N, K, U_{\max})$ such that for every $\theta, \theta'$ and
every $t$ with $\varepsilon_t \leq 1$,
$|\tilde Q^{(t)}(\theta;\theta') - Q_{\mathrm{true}}(\theta;\theta')|
\leq C_Q\,\varepsilon_t$.
\end{lemma}

\begin{proof}
Add and subtract
$\sum_m |\mathcal{C}_m| \sum_k q^{(t)}_{mk}\bigl[\log r_{mk}
- u_{\mathrm{true}}(\boldsymbol{\alpha}_k, \mathbf{r}_k)\bigr]$.
The first piece, the surrogate-error term, is
$\sum_m |\mathcal{C}_m| \sum_k q^{(t)}_{mk}
(u_{\mathrm{true}} - u^{(t)})$ and is bounded by $N\varepsilon_t$.
The remainder splits as
$\sum_m |\mathcal{C}_m| \sum_k (q^{(t)}_{mk} - w_{mk})(-u_{\mathrm{true}})$
and
$\sum_m |\mathcal{C}_m| \sum_k (q^{(t)}_{mk} - w_{mk})\log r_{mk}$.
The first is bounded by $3NK\,U_{\max}\,\varepsilon_t$ via
Lemma~\ref{lem:resp} (because $|u_{\mathrm{true}}| \leq U_{\max}$).
For the second, using $|x\log x|\leq 1/e$ on $[0,1]$ together with
the softmax-stability bound of Lemma~\ref{lem:softmax}, one obtains a
bound of $3N\varepsilon_t/e$. Summing the three contributions gives
the claim with $C_Q = N(1 + 3K\,U_{\max} + 3/e)$.
\end{proof}

\subsection{Main Proof}

\begin{proposition}[Quasi-monotonicity of $L_{\mathrm{true}}$]
\label{prop:quasi-mono}
Under Assumptions~\ref{ass:bounded}--\ref{ass:gem}, for every $t$,
\begin{equation}
\label{eq:quasi-mono}
L_{\mathrm{true}}(\theta^{(t+1)}) \;\geq\;
L_{\mathrm{true}}(\theta^{(t)}) \;-\; 2N\,\varepsilon_t.
\end{equation}
\end{proposition}

\begin{proof}
We chain five inequalities. First,
$L_{\mathrm{true}}(\theta^{(t+1)}) \geq
\mathcal{F}_{u_{\mathrm{true}}}(q^{(t)}, \theta^{(t+1)})$ since
$L_u(\theta) = \max_q \mathcal{F}_u(q,\theta)$. Second, by
Lemma~\ref{lem:elbo-dev},
$\mathcal{F}_{u_{\mathrm{true}}}(q^{(t)}, \theta^{(t+1)})
\geq \mathcal{F}_{u^{(t)}}(q^{(t)}, \theta^{(t+1)}) - N\varepsilon_t$.
Third, the GEM Assumption~\ref{ass:gem} yields
$\mathcal{F}_{u^{(t)}}(q^{(t)}, \theta^{(t+1)})
\geq \mathcal{F}_{u^{(t)}}(q^{(t)}, \theta^{(t)})$. Fourth, the
E-step makes the latter equal to $L^{(t)}(\theta^{(t)})$. Fifth,
Corollary~\ref{cor:lik-dev} gives $L^{(t)}(\theta^{(t)}) \geq
L_{\mathrm{true}}(\theta^{(t)}) - N\varepsilon_t$. Combining the
five steps yields~\eqref{eq:quasi-mono}.
\end{proof}

\textbf{Part (a): convergence of the likelihood.}
Set $L_{\max} = \sup_{\theta\in\Theta} L_{\mathrm{true}}(\theta)$,
which is finite because $\Theta$ is compact (Assumption~\ref{ass:cont}
and finiteness of $\mathcal{A}$). Define
$V_t = L_{\max} - L_{\mathrm{true}}(\theta^{(t)}) \geq 0$. From
Proposition~\ref{prop:quasi-mono},
\begin{equation*}
V_{t+1} \;\leq\; V_t \;+\; 2N\varepsilon_t,
\qquad \sum_{t=0}^{\infty} 2N\varepsilon_t < \infty
\;\;\text{by Assumption~\ref{ass:summable}.}
\end{equation*}
The Robbins--Siegmund theorem~\cite{robbins1971convergence} (with
$a_t = 0$, $b_t = 2N\varepsilon_t$) implies that $\{V_t\}$ converges,
hence $\{L_{\mathrm{true}}(\theta^{(t)})\}$ converges. \qed

\textbf{Part (b): stationarity of limit points.}

\textbf{Step 1: existence.} $\Theta$ is compact, so
$\{\theta^{(t)}\}$ has at least one limit point
$\theta^{\star} = \lim_{j\to\infty} \theta^{(t_j)}$.

\textbf{Step 2: contradiction setup.} Suppose $\theta^{\star}$ is not
stationary: there is $\theta'\in\Theta$ and $\delta > 0$ with
$Q_{\mathrm{true}}(\theta';\theta^{\star}) -
Q_{\mathrm{true}}(\theta^{\star};\theta^{\star}) = \delta > 0$.

\textbf{Step 3: continuity.} The map $\theta'\mapsto Q_{\mathrm{true}}(\theta;\theta')$
is continuous on $\Theta$: $w_{mk}(\theta')$ is continuous because
$u_{\mathrm{true}}$ is continuous in $\mathbf{r}$
(Assumption~\ref{ass:cont}) and softmax is continuous; and the
product $w_{mk}\log r_{mk}$ extends continuously to $r_{mk} = 0$
because $w_{mk} = O(r_{mk})$ and $r\mapsto r\log r$ is continuous at
$0$. Thus, for $j$ large enough,
$Q_{\mathrm{true}}(\theta';\theta^{(t_j)}) -
Q_{\mathrm{true}}(\theta^{(t_j)};\theta^{(t_j)}) \geq \delta/2$.

\textbf{Step 4: transfer to $\tilde Q$.} By Lemma~\ref{lem:qdev},
\begin{equation*}
\tilde Q^{(t_j)}(\theta';\theta^{(t_j)}) -
\tilde Q^{(t_j)}(\theta^{(t_j)};\theta^{(t_j)}) \;\geq\;
\delta/2 - 2C_Q\,\varepsilon_{t_j} \;\geq\; \delta/4
\end{equation*}
for $j$ large enough, since $\varepsilon_{t_j}\to 0$.

\textbf{Step 5: GEM ascent.} Because the M-step searches over a
candidate set that contains $\theta'$ (by finiteness of $\mathcal{A}$
and the closure properties of the $\mathbf{R}$-update),
$\tilde Q^{(t_j)}(\theta^{(t_j+1)};\theta^{(t_j)}) \geq
\tilde Q^{(t_j)}(\theta';\theta^{(t_j)})$. Translating back through
Lemma~\ref{lem:qdev} once more,
\begin{equation*}
\Delta_{t_j} := Q_{\mathrm{true}}(\theta^{(t_j+1)};\theta^{(t_j)}) -
Q_{\mathrm{true}}(\theta^{(t_j)};\theta^{(t_j)}) \;\geq\;
\delta/4 - 2C_Q\,\varepsilon_{t_j} \;\geq\; \delta/8
\end{equation*}
for $j$ large enough.

\textbf{Step 6: summability contradiction.} The Gibbs inequality
$L_{\mathrm{true}}(\theta^{(t+1)}) - L_{\mathrm{true}}(\theta^{(t)})
\geq \Delta_t$ together with Proposition~\ref{prop:quasi-mono}
implies $\Delta_t \geq -2N\varepsilon_t$. Hence the non-negative
sequence $\Delta_t + 2N\varepsilon_t$ satisfies
$\sum_t (\Delta_t + 2N\varepsilon_t) < \infty$ by
Robbins--Siegmund applied to $V_t$, and consequently
$\sum_t \Delta_t < \infty$. But Step~5 gives $\Delta_{t_j} \geq
\delta/8$ for infinitely many $j$, yielding $\sum_t \Delta_t = +\infty$
--- a contradiction. Therefore $\theta^{\star}$ is stationary. \qed

\subsection{Convergence Rate}

A by-product of the proof of Theorem~\ref{thm:sgem}(a) is the
explicit rate
\begin{equation}
L_{\mathrm{true}}(\theta^{\star}) - L_{\mathrm{true}}(\theta^{(T)})
\;\leq\; 2N \sum_{t = T}^{\infty} \varepsilon_t.
\end{equation}
If $\varepsilon_t = O(t^{-p})$ with $p > 1$, this gives
$L_{\mathrm{true}}(\theta^{(T)}) \to L_{\mathrm{true}}(\theta^{\star})$
at rate $O(T^{1 - p})$, matching the rate observed for stochastic-EM
analyses~\cite{cappe2009online}.


\label{sec:appendix}

\end{document}
